\documentclass{article}
\usepackage{mathnotes}

\title{Outer Measure}
\slug{outer-measure}
\description{Outer measure on the real line.}

\begin{document}

\section{Outer Measure}

Our intention is to come up with an \@{extension} of the \@{length} of \@{intervals} that works on more complicated \@{domains} than the \@{subintervals} that \@{Riemann integration} uses, for example, \@{unions} of \@{subsets} of the \@{reals}.

\begin{definition}[length of open interval]
\notation{\length}{\ell}

	The \textbf{length $\length(I)$ of an \@{open interval}} $I$ is defined by

	\[ \begin{cases}
		b - a & \text{ if } I = (a,b) \text{ for some } a,b \in \reals \text{ with } a < b, \\
		0 & \text{ if } I = \emptyset, \\
		\infty & \text{ if } I = (-\infty,a) \text{ or } I = (a, \infty) \text{ for some } a \in \reals, \\
		\infty & \text{ if } I = (-\infty,\infty) \\
	\end{cases} \]
\end{definition}

\begin{definition}[outer measure]
	The \term{outer measure} $|A|$ of a \@{set} $A \subset \reals$ is defined by

	\[ |A| = \inf \left \{ \sum_{k=1}^{\infty} \length(I_k) : I_1, I_2, \dots \text{ are open intervals such that } A \subset \bigcup_{k=1}^\infty I_k \right \}. \]
\end{definition}

\subsection{Properties of Outer Measure}\label{outer-measure-properties}
TODO: I have proofs for these written down in paper notes, transcribe them.

\begin{theorem}[Countable Sets Have Outer Measure 0]
	Every \@{countable} \@{subset} of $\reals$ has \@{outer measure} 0.
\end{theorem}
\begin{corollary}
	From this and \@{rationals-are-countable}, we have that $|\rationals| = 0.$
\end{corollary}

\begin{theorem}[Outer Measure Preserves Order]
	Suppose $A$ and $B$ are \@{subsets} of $\reals$ with $A \subset B.$ Then $|A| \leq |B|.$
\end{theorem}

\begin{definition}[translation]
	If $t \in \reals$ and $A \subset \reals,$ then the \term{translation} $t + A$ is defined by

	\[ t + A = \{ t + a : a \in A \}. \]
\end{definition}

\begin{theorem}[Outer measure is translation invariant]
Suppose $t \in \reals$ and $A \subset \reals.$ Then $|t + A| = |A|.$
\end{theorem}

\begin{theorem}[Countable subadditivity of outer measure]
	Suppose $A_1, A_2, \dots$ is a \@{sequence} of \@{subsets} of $\reals.$ Then

	\[ \left | \bigcup_{k=1}^{\infty} A_k \right | \leq \sum_{k=1}^{\infty} |A_k|.  \]
\end{theorem}

\begin{theorem}[Outer measure of a closed interval]
	Suppose $a,b \in \reals$ with $a < b,$ Then $|[a,b]| = b - a.$
\end{theorem}
\begin{note}
	The proof for this theorem uses \@{heine-borel}.
\end{note}

\begin{theorem}[Nontrivial intervals are uncountable]
	Every \@{interval} in $\reals$ that contains at least two \@{distinct} \@{elements} is \@{uncountable}.
\end{theorem}


Now, the problem with \@{outer measure}, or really, the problem with $\reals.$

\begin{theorem}[Nonadditivity of outer measure]
	There exist \@{disjoint} \@{subsets} $A$ and $B$ or $\reals$ such that

	\[ |A \cup B| \neq |A| + |B|. \]
\end{theorem}
\begin{note}
	The proof of this theorem uses \href{https://en.wikipedia.org/wiki/Vitali_set}{Vitali Sets}. 
\end{note}


\end{document}
